\documentclass[11pt]{article}
\usepackage[margin=1.1in,top=1.0in,bottom=1.0in]{geometry}
\usepackage{amsmath,amssymb}
\usepackage{booktabs}
\usepackage{hyperref}
\usepackage{xcolor}
\usepackage{enumitem}
\usepackage{microtype}

\hypersetup{colorlinks=true,linkcolor=blue,citecolor=blue,urlcolor=blue}

\newcommand{\PDG}{\textsc{pdg}}
\newcommand{\RA}{\textsc{ra}}

\begin{document}

\begin{center}
  {\large\bfseries Relational Actualism: A Research Programme Summary}\\[6pt]
  {\normalsize Joshua F.\ Sandeman\quad\textbar\quad
   Independent Researcher, Salem, Oregon\quad\textbar\quad
   \href{mailto:sansuikyo@gmail.com}{sansuikyo@gmail.com}}\\[4pt]
  {\small
   Zenodo community: \href{https://zenodo.org/communities/relational-actualism}{zenodo.org/communities/relational-actualism}\quad\textbar\quad
   GitHub: \href{https://github.com/jsandeman/Relational-Actualism}{github.com/jsandeman/Relational-Actualism}\\
   Website: \href{https://relationalactualism.org}{relationalactualism.org}\quad\textbar\quad
   April 2026}
\end{center}

\bigskip
\hrule
\bigskip

\noindent
\textbf{The single commitment.}
Relational Actualism (RA) begins from one ontological claim: on-shell actualization
events — irreversible physical interactions satisfying $\Delta S(\rho\|\sigma_0) > 0$
that write permanent vertices into a growing causal Directed Acyclic Graph —
are the primitive constituents of physical reality.
Time is not a background parameter through which events are located;
time is the process of the causal graph growing.

From this single commitment, and without introducing additional fields,
parameters, or new physics, the framework derives the results summarized below.

\bigskip

% ─────────────────────────────────────────────────────────────────────────────
\noindent\textbf{Key quantitative results.}

\medskip
\noindent\textit{Standard Model coupling constants and particle spectrum}
(Papers RASM, RATM):

\begin{center}
\small
\begin{tabular}{lccc}
\toprule
Quantity & RA prediction & \PDG\ value & Agreement \\
\midrule
$\alpha_{\mathrm{EM}}^{-1}$ (leading order) & $144 - 7 = 137.000$ & $137.036$ & $0.026\%$ \\
$\alpha_{\mathrm{EM}}^{-1}$ (BDG fixed-point) & $136.845$ & $137.036$ & $0.14\%$ \\
$\alpha_s^{\mathrm{UV}}$ (BDG RG UV fixed point) & $1/\sqrt{72} = 0.11785$ & $0.11800 \pm 0.0009$ & $0.13\%$ \\
$\alpha_s^{\mathrm{IR}}$ (BDG RG IR fixed point) & $1/3 = 0.333$ & $\sim\!0.3$ at $1\;\mathrm{GeV}$ & consistent \\
$m_p/\Lambda_{\mathrm{QCD}}$ & $\sqrt{c_4} = 2\sqrt{2} = 2.828$ & $2.826$ & $0.08\%$ \\
$\sin^2\theta_W|_{\mathrm{GUT}}$ & $3/8 = 0.375$ & $0.375$ (SU(5)) & exact \\
SM particle families & 5 BDG topology classes & 5 families & exact bijection \\
\bottomrule
\end{tabular}
\end{center}

\medskip
\noindent All five Standard Model gauge coupling constants and the
complete particle spectrum emerge from the five integers
$(1,-1,9,-16,8)$ that uniquely specify the Benincasa--Dowker--Glaser
causal set action in $d=4$ [BDG2010, Glaser2014].
The BDG renormalization group has two exact fixed points
($\alpha_s^{\mathrm{IR}} = 1/3$ at confinement,
$\alpha_s^{\mathrm{UV}} = 1/\sqrt{72}$ at asymptotic freedom) derived
from these integers alone.

\medskip
\noindent\textit{Quantum foundations} (Papers RAQM, RAEB):
The actualization criterion $\Delta S(\rho\|\sigma_0) > 0$
provides an observer-independent, basis-independent, non-perturbative
criterion for wave function collapse.
It dissolves the measurement problem (collapse is a physical event,
not an observer-relative postulate), resolves the Unruh paradox
(both inertial and Rindler observers agree $\Delta S = 0$ in the vacuum),
and eliminates the Heisenberg cut (classicality is a gradient in
actualization density, not a boundary).
Three experimental predictions discriminate this framework from all current
interpretations: a finite Lieb-Robinson propagation velocity for the
actualization front; an environment-dependent transition duration
$\tau_{\mathrm{act}} \sim \hbar/\Delta E_{\mathrm{env}}$; and a sharp,
measurable reversibility threshold $\Delta S^*$.

\medskip
\noindent\textit{Cosmology and dark matter} (Papers RAGC, RADM):
The cosmological constant problem does not arise in RA: the vacuum does
not gravitate because off-shell processes produce no actualization events.
The metric sourcing equation $\nabla^2\Phi = 4\pi G[\rho_A + \rho_\lambda]$
(where $\rho_\lambda \propto \nabla^2(\ln\lambda)$) produces flat galactic
rotation curves without dark matter, explains the $4$--$6\sigma$ Hubble
tension as a density-dependent expansion effect, and categorically prohibits
WIMPs from gravitating.
Falsifiable predictions: $H_{\mathrm{Eridanus}} \approx 75.9\;\mathrm{km\,s^{-1}\,Mpc^{-1}}$
(testable with current void survey data) and a gravitational lensing
Einstein radius $\theta_E \approx 0.47\;\mathrm{arcsec}$ for a Milky
Way analog (distinguishable from $\Lambda$CDM at $0.9$--$1.4\;\mathrm{arcsec}$).

\medskip
\noindent\textit{Emergence and complexity} (Papers RAHC, RACI, RACF):
Biological assembly corresponds directly to actualization events.
The Causal Firewall is the $\mu = 1$ Erd\H{o}s-R\'{e}nyi percolation
transition: the same critical actualization density governs QCD
confinement--deconfinement, the onset of flat galactic rotation curves,
fault-tolerance thresholds in quantum error correction, and the onset
of stable biological complexity.
The framework predicts substrate-independent biosignatures correlated with
environments sustaining $\mu \geq 1$.

\bigskip

% ─────────────────────────────────────────────────────────────────────────────
\noindent\textbf{Lean~4 formal verification.}
The discrete graph-theoretic and algebraic core of the framework is
formally verified in Lean~4/Mathlib.
Key machine-checked results include: the BDG particle topology closure
(124 extension cases; no new SM family can emerge), colour confinement
as a theorem ($L_q = 4$, $L_g = 3$), the $\alpha_{\mathrm{EM}}^{-1} = 137$
integer identity, frame independence
($S(U\rho U^\dagger\|\sigma_0) = S(\rho\|\sigma_0)$),
Rindler stationarity ($\Delta S = 0$ under modular flow),
and the Local Ledger Condition at causal severances.
Three proof files: \texttt{RA\_D1\_Proofs.lean} (73 theorems),
\texttt{RA\_Alpha\_EM\_Proof.lean} (20 theorems), and
\texttt{RA\_AQFT\_Proofs\_v10.lean} — all compiling in Lean~4.29.0.
All files available at the GitHub repository.

\bigskip

% ─────────────────────────────────────────────────────────────────────────────
\noindent\textbf{Submission and publication status (April 2026).}

\begin{center}
\small
\begin{tabular}{p{4.4cm}lll}
\toprule
Paper & Journal & Status & Zenodo DOI \\
\midrule
RAQM (quantum mechanics) & \textit{Foundations of Physics} & Under review & 10.5281/zenodo.19174380 \\
RACL (EFE uniqueness) & \textit{Class.\ Quantum Grav.} & Under review & 10.5281/zenodo.19270020 \\
RATM (particle topology) & \textit{Physical Review D} & Under review & 10.5281/zenodo.19265651 \\
RACF (Causal Firewall) & \textit{Int.\ J.\ Astrobiology} & Under review & 10.5281/zenodo.19319374 \\
\midrule
RAGC, RAEB, RADM, RAQI & \textit{PRD, CQG, JCAP, PRA} & Ready & see Zenodo community \\
RAHC, RACI & \textit{PRE / Entropy} & Ready & see Zenodo community \\
RASM & \textit{Annals of Physics} & Preparation & 10.5281/zenodo.19229335 \\
\bottomrule
\end{tabular}
\end{center}

\medskip
\noindent The complete twelve-paper suite is
available at the Zenodo community:
\href{https://zenodo.org/communities/relational-actualism}{zenodo.org/communities/relational-actualism}.

\bigskip

% ─────────────────────────────────────────────────────────────────────────────
\noindent\textbf{Open problems and collaboration targets.}

The framework has four precisely identified mathematical gaps:

\begin{enumerate}[itemsep=3pt,leftmargin=*]

\item \textbf{Bianchi compatibility.} The RA field equation
$G_{\mu\nu} = 8\pi G(T^{(A)}_{\mu\nu} + \Theta^{(\lambda)}_{\mu\nu})$
must satisfy $\nabla_\mu G^{\mu\nu} \equiv 0$ covariantly.
The weak-field limit holds by construction; the full covariant proof
requires deriving the coupling $\xi$ from the microscopic graph dynamics
via Tomita-Takesaki modular theory.
\textit{Collaboration target:} algebraic QFT, type III$_1$ von Neumann algebras.

\item \textbf{Amplitude locality (discrete intrinsic proof).}
The causal invariance theorem is proved in the perturbative regime
given amplitude locality as an axiom (physically justified by QFT
microcausality).
The discrete intrinsic formulation — proving Bell causality for the
graph rewriting system — is the primary open problem of the Lean
verification programme.
\textit{Collaboration target:} causal set theory (Dowker, Henson, Sorkin group).

\item \textbf{Continuum limit.}
The discrete Local Ledger Condition (Lean-verified) must be shown to
converge to the contracted Bianchi identity $\nabla_\mu G^{\mu\nu}
\equiv 0$ without breaking Lorentz invariance in the
$\lambda \to \infty$ limit.
\textit{Collaboration target:} causal dynamical triangulations (Loll/Radboud),
asymptotic safety.

\item \textbf{BDG density-to-energy scale mapping.}
The BDG renormalization group's dimensionless density parameter $\mu$
must be mapped to a physical energy scale via the RA continuum limit.
This would convert the UV fixed-point agreement with $\alpha_s(m_Z)$
from a numerical coincidence into a first-principles derivation.
\textit{Collaboration target:} lattice QCD, causal set dynamics.

\end{enumerate}

\bigskip

% ─────────────────────────────────────────────────────────────────────────────
\noindent\textbf{An invitation.}
The programme is an explicit call for engagement, not a completed theory.
Every open problem above is a well-posed mathematical target.
The suite's predictions are specific, falsifiable, and disciplinarily
diverse: they span particle physics, quantum gravity, quantum information,
cosmology, and astrobiology.
If any of the claims in this summary intersect your research,
I would welcome correspondence.

\bigskip
\noindent The full suite, Lean proof files, and Python computational
scripts are publicly available. This is an independent research programme
— the work is what it is, and the community will judge it.

\begin{flushright}
Joshua F.\ Sandeman\\
Independent Researcher, Salem, Oregon\\
\href{mailto:sansuikyo@gmail.com}{sansuikyo@gmail.com}\\
April 2026
\end{flushright}

\bigskip
\hrule
\bigskip

\begingroup
\small
\noindent\textbf{Selected references}

\noindent
[BDG2010] D.\,M.\,T.\ Benincasa and F.\ Dowker,
\textit{The scalar curvature of a causal set},
Phys.\ Rev.\ Lett.\ \textbf{104}, 181301 (2010).

\noindent
[Glaser2014] L.\ Glaser,
\textit{A closed form expression for the causal set d'Alembertian},
Class.\ Quantum Grav.\ \textbf{31}, 095007 (2014).

\noindent
[Sorkin1994] R.\,D.\ Sorkin,
\textit{Quantum mechanics as quantum measure theory},
Mod.\ Phys.\ Lett.\ A \textbf{9}, 3119 (1994).

\noindent
[PDG2024] R.\,L.\ Workman et al.\ (Particle Data Group),
Prog.\ Theor.\ Exp.\ Phys.\ \textbf{2022}, 083C01; updated 2024.

\noindent
[Lean4] L.\ de Moura and S.\ Ullrich, Lean~4 theorem prover,
Proc.\ CADE-28, LNCS \textbf{12699}, 625 (2021).
\endgroup

\end{document}
